\section{Related Work}
\label{sec:related_work}

The field of model merging has experienced a surge of interest as a cost-effective alternative to multi-task fine-tuning or joint training. In this section, we review the foundational paradigms of static model merging, dynamic merging and routing, and the emergence of regularization techniques in this domain.

\subsection{Static Model Merging}
Static model merging combines multiple specialized models fine-tuned from a common progenitor into a single set of weights, using fixed coefficients across all inputs. Early approaches, such as weight averaging \cite{wortsman2022robust}, demonstrated that simple linear interpolation of models fine-tuned on the same task with different hyperparameters or seeds can boost generalization. This was extended to multi-task settings by \citet{ilharco2022editing} through \emph{Task Arithmetic}, which defines a task vector as the difference between the specialized model and the base model weights. Task vectors can be scaled and added to the base model to inject task-specific capabilities. 

However, scaling and adding multiple task vectors often leads to severe interference in parameter space, particularly for highly overlapping or conflicting parameters. To mitigate this, several advanced static merging methods have been developed. \citet{tiesmerging} introduced \emph{TIES-Merging} (Trim, Elect, and Sign), which addresses sign conflicts and redundant parameter values by trimming low-magnitude values, electing a dominant sign, and merging only aligned parameters. Similarly, \citet{yu2024dare} proposed \emph{DARE} (Drop and Rescale), which randomly drops a high percentage of task vector parameters and rescales the remainder to maintain the expectation of the model weights. While these methods achieve remarkable performance in static settings, they are fundamentally unable to adapt to different task balances or input domains dynamically at inference time.

\subsection{Dynamic Model Merging and Routing}
To overcome the rigid nature of static merging, dynamic model merging introduces input-dependent or batch-dependent coefficients. \citet{adamerging} proposed \emph{AdaMerging}, which learns task-specific merging coefficients at a layer-wise or block-wise granularity. While AdaMerging learns these coefficients during an offline calibration phase, they remain static during inference. 

To achieve true input-adaptive merging, recent works have proposed dynamic routing heads that compute merging coefficients on the fly for each input sample or batch. For instance, \citet{adamerging} also introduced an active routing mechanism where a lightweight classifier predicts the task label of the input and uses the prediction confidence to scale the task expert coefficients. However, this relies on explicit task classification and can fail when inputs are ambiguous or represent a mixture of domains. More advanced dynamic routers, such as the Linear Router or Bounded Linear Router (BL-Router) \cite{adamerging}, map intermediate activation representations directly to merging coefficients. 

While dynamic routing heads offer exceptional flexibility, they introduce a significant optimization challenge: the routing head must be calibrated on extremely small datasets (e.g., a few dozen samples) to remain practical. Under these low-data regimes, unregularized optimization of the routing head can lead to catastrophic overfitting or representational collapse. This is particularly pronounced on high-conflict tasks (such as SVHN when mixed with MNIST or CIFAR-10), where the routing head fails to allocate appropriate pathways for the difficult task.

\subsection{Wave-inspired Merging and Regularization}
In response to the limitations of classical linear routers, some recent works have proposed complex physical or mathematical metaphors to handle task conflicts. For example, \emph{Quantum Wavefunction Superposition Merging} (QWS-Merge) models task Experts as wavefunctions in a complex-valued Hilbert space, resolving conflicts through simulated phase interference and spherical projections. While these methods demonstrate improved performance, they often require complex, non-standard optimization routines and introduce heavy mathematical overhead.

In contrast, our work focuses on systematic empirical validation. We argue that the success of complex wave-inspired merging is primarily due to their implicit regularization properties, rather than the intrinsic correctness of their physical metaphors. By systematically isolating and deconstructing these failures, we show that simple, classical sigmoidal routing heads, when equipped with appropriate prior regularization (such as our proposed TCPR), can achieve superior performance without any of the conceptual or computational overhead of wave-inspired models.
